% =====================================================================
% §1.5 — Diagrama de flujo del proceso productivo
% =====================================================================
\section{Flujo del proceso productivo}
\label{sec:flujo}

El flujo abarca desde la preparación e inoculación del sustrato hasta
la distribución del producto final, fresco o deshidratado.

% ---------------------------------------------------------------------
\subsection{Etapas y duración orientativa}
\label{sec:flujo:etapas}

\begin{table}[H]
  \centering
  \caption{Etapas del proceso productivo. Tiempos orientativos para
    \paramEspecieA{}, ajustables por especie (\S\ref{sec:especies}).}
  \label{tab:flujo}
  \begin{tabularx}{\linewidth}{lXl}
    \toprule
    \textbf{Etapa} & \textbf{Acción} & \textbf{Duración} \\
    \midrule
    1 & Hidratación del sustrato y mezcla con suplementos.            & 12--24\,h \\
    2 & Embolsado en bolsa autoclavable con parche filtro.            & 0{,}5\,h/lote \\
    3 & Esterilización en olla de presión (\S\ref{sec:preparacion}).  & \paramTiempoEsterSustrato\,min + enfriado 12\,h \\
    4 & Inoculación con grain spawn en SAB.                           & 0{,}3\,h/bolsa \\
    5 & Incubación en oscuridad, T ambiente.                          & 10--21\,días \\
    6 & Inducción de fructificación (al contenedor sellado).          & 0\,h (transición) \\
    7 & Pinning + maduración.                                         & 5--14\,días \\
    8 & Cosecha.                                                      & 0{,}5\,h/contenedor \\
    9a& Empaque fresco $\to$ refrigeración 2--4\,\si{\celsius}.       & inmediato \\
    9b& Deshidratado (\S\ref{sec:postcosecha}) $\to$ empaque hermético.& 6--14\,h \\
    \bottomrule
  \end{tabularx}
\end{table}

\subsection{Diagrama general}
\label{sec:flujo:diagrama}

El diagrama renderizable se encuentra en
\texttt{diagramas/flujo-productivo.mmd} (mermaid). Para incluirlo en el
PDF se exporta a PNG y se inserta:

\begin{verbatim}
mmdc -i diagramas/flujo-productivo.mmd -o diagramas/flujo-productivo.png
\end{verbatim}

\subsection{Puntos de registro hacia la base de datos (Tópico~7)}
\label{sec:flujo:registro}

Cada etapa con \emph{flecha hacia DB} en el diagrama
\texttt{flujo-productivo.mmd} corresponde a un registro persistente:

\begin{table}[H]
  \centering
  \caption{Puntos de registro a la base de datos unificada.}
  \label{tab:registro}
  \begin{tabularx}{\linewidth}{lXl}
    \toprule
    \textbf{Etapa} & \textbf{Variable registrada} & \textbf{Tabla destino} \\
    \midrule
    1 & Lote, especie, sustrato, masa hidratada                   & \texttt{lotes} \\
    3 & T y t de esterilización, presión registrada               & \texttt{lotes} \\
    4 & Fecha inoculación, jeringa/spawn usado                    & \texttt{lotes} \\
    5 & Detección colonización completa (inspección)              & \texttt{lotes} \\
    6 & Fecha apertura, contenedor destino                        & \texttt{lotes} \\
    7 & Telemetría continua HR, T, CO$_2$, lux (30\,s)            & \texttt{mediciones} \\
    7 & Imagen cada 15\,min                                       & \texttt{imagenes} \\
    8 & Masa fresca cosechada, BE\,\% calculado                   & \texttt{cosechas} \\
    9a/9b & Decisión fresco/seco; masa fresca destinada a secado  & \texttt{cosechas} \\
    9b & Masa seca final, factor conversión, HR residual          & \texttt{cosechas} \\
    \bottomrule
  \end{tabularx}
\end{table}

\subsection{Criterios de aceptación}
\label{sec:flujo:aceptacion}

\begin{itemize}[leftmargin=*,itemsep=0pt]
  \item El diagrama identifica todo punto de registro y lo mapea a una
        tabla del schema (Tópico~7).
  \item Cada transición tiene una condición de paso explícita
        (colonización completa, primordios visibles, sombrero
        aplanado, HR residual $<$ objetivo).
  \item Los tiempos de tabla~\ref{tab:flujo} se actualizan con datos
        propios tras los primeros 3 lotes.
\end{itemize}
